CUDA Fortranで高次精度の差分を計算する際はshared memoryの使用が推奨される．Shared memoryを使う場合はshared属性をつけて配列を宣言する．Shared memoryの形状にはparameter属性を付与してコンパイル時に形状が決まるようにするとよい．Shared memoryを動的に確保することもできるが，静的に確保した方が速い．

6点stencilの差分計算において，global memoryからshared memoryに素朴にコピーしてみる．Qをglobal memory上の配列，rhoをshared memory上の配列とする．

\begin{lstlisting}[language=Fortran]
rho(it-2:it+3,jt,kt) = Q(1,i-2:i+3,j,k)
\end{lstlisting}

\noindent
このようにすると複数のスレッドが同一のshared memory上の番地にアクセスするため”待ち”が発生する．これをBank conflictという．Bank conflictを解消するには以下のように書くとよい．なお，syncthreadsを呼び忘れるとスレッド間で同期がとられないため以降の差分計算が失敗する．

\begin{lstlisting}[language=Fortran]
it = threadIdx%x
i_base = (blockIdx%x-1)*blockDim%x
do ii = it-2, threadsE%x+3, blockDim%x
  i = i_base + ii
  if (i >= 1 .and. i <= nx) then
    rho(ii,jt,kt) = Q(1,i,j,k)
  endif
enddo
call syncthreads()
\end{lstlisting}

\noindent
CUDA Fortranで究極的な実行速度を達成するには一次元配列を用いるのが望ましい．以下に一次元のshared memoryを用いた例を示す．OUxSBLIでは以下のようにglobal memoryからshared memoryへのコピーを行っている．（2026年3月時点）

\begin{lstlisting}[language=Fortran]
it = threadIdx%x
jt = threadIdx%y
kt = threadIdx%z
j  = (blockIdx%y-1)*blockDim%y + jt + 1
k  = (blockIdx%z-1)*blockDim%z + kt + 1
i_base = (blockIdx%x-1)*blockDim%x
offset_yz = (jt-1) * sx + (kt-1) * sx * sy
do ii = it-2, threadsE%x+3, blockDim%x
  i = i_base + ii
  if (i >= 1 .and. i <= nx .and. j >= 1 .and. j <= ny .and. k >= 1 .and. k <= nz) then
    idx = ii + offset_yz
    rho(idx) = Q(1,i,j,k)
      u(idx) = Q(2,i,j,k)
      v(idx) = Q(3,i,j,k)
      w(idx) = Q(4,i,j,k)
      p(idx) = Q(5,i,j,k)
    tmp(idx) =   T(i,j,k)
  endif
enddo
call syncthreads()
\end{lstlisting}

CUDA Fortranでは連続アクセスでない部分配列をsubroutineに渡すことができない．したがって，y方向＆z方向の差分を計算する際に問題が生じる．OUxSBLIではy方向＆z方向の差分を計算する際はshared memory上で差分方向に連続アクセスになるように転置するという方式を採用している．単純な中心差分では転置のコストで全体が遅くなる可能性が考えられるが，高解像度風上差分を用いる場合は転置のコストよりも連続アクセスであることによるその後の処理の恩恵の方が大きいと考えている．

\begin{lstlisting}[language=Fortran]
it = threadIdx%x
jt = threadIdx%y
kt = threadIdx%z
i  = (blockIdx%x-1)*blockDim%x + it + 1
k  = (blockIdx%z-1)*blockDim%z + kt + 1
j_base = (blockIdx%y-1)*blockDim%y
offset_xz = (it-1) * sy + (kt-1) * sy * sx
do jj = jt-2, threadsF%y+3, blockDim%y
  j = j_base + jj
  if (i >= 1 .and. i <= nx .and. j >= 1 .and. j <= ny .and. k >= 1 .and. k <= nz) then
    idx = jj + offset_xz
    rho(idx) = Q(1,i,j,k)
      u(idx) = Q(2,i,j,k)
      v(idx) = Q(3,i,j,k)
      w(idx) = Q(4,i,j,k)
      p(idx) = Q(5,i,j,k)
    tmp(idx) =   T(i,j,k)
  endif
enddo
call syncthreads()
\end{lstlisting}
